\documentclass[11pt,a4paper]{article}

\usepackage[utf8]{inputenc}
\usepackage[T1]{fontenc}
\usepackage{amsmath,amssymb}
\usepackage{booktabs}
\usepackage{hyperref}
\usepackage{xcolor}
\usepackage{colortbl}
\usepackage[margin=2.5cm]{geometry}
\usepackage{caption}
\usepackage{subcaption}
\usepackage{array}

\definecolor{bestcell}{rgb}{0.85,1.0,0.85}
\definecolor{badcell}{rgb}{1.0,0.88,0.85}
\definecolor{tlcell}{rgb}{0.95,0.95,0.95}
\definecolor{headercolor}{rgb}{0.22,0.39,0.60}

\newcommand{\best}[1]{\cellcolor{bestcell}\textbf{#1}}
\newcommand{\bad}[1]{\cellcolor{badcell}{#1}}

\title{\textbf{Evaluation of LP Relaxation Solver Strategies in CBC}\\
       \large Analysis of \texttt{guess()} Parameter Configurations}
\author{CBC Development}
\date{March 2026}

\begin{document}
\maketitle
\tableofcontents
\newpage

% ─────────────────────────────────────────────────────────────────────────────
\section{Overview}

The CBC solver's \texttt{-guess} option (implemented in
\texttt{ClpSimplexOther::guess()}) selects one of four LP relaxation
configurations based on instance characteristics (integrality of variables and
statistics of objective coefficients). The goal of this experiment is to evaluate
each configuration individually on a representative benchmark set in order to:
\begin{itemize}
  \item identify correctness problems (wrong objective, primal infeasibilities);
  \item identify reliability problems (time-limit kills, no solution produced);
  \item determine which configurations have a genuine performance advantage over
        the others, and whether any can be removed from the selection pool.
\end{itemize}

All runs used the ASan+debug build (\texttt{-O1 -g -fsanitize=address}) installed at
\texttt{\textasciitilde/prog/cbc-dbg/bin/cbc}, to additionally detect memory errors
and assertion failures.

% ─────────────────────────────────────────────────────────────────────────────
\section{Experimental Setup}

\subsection{Benchmark Set}
The benchmark consists of \textbf{1\,317 MIP instances} from the
\texttt{\textasciitilde/inst/super/} collection. Reference LP relaxation
objectives were obtained from Gurobi
(\texttt{lp\_relaxation\_results\_gurobi.tsv}).

\subsection{Solver Configurations}
The four configurations correspond exactly to the branches inside
\texttt{guess()} and are selected by two binary conditions: whether all
variables with $\ell < u$ are integer (\texttt{allInt}), and whether the median
or average absolute objective coefficient exceeds a threshold.

\begin{table}[h]
\centering
\caption{The four LP solver configurations tested.}
\begin{tabular}{llp{7.5cm}}
\toprule
\textbf{Name} & \textbf{Condition} & \textbf{CBC options} \\
\midrule
\texttt{idiot30} & \texttt{allInt}, avg $\le 0.0086$ &
  \texttt{-idiot 30 -pertvalue -1483 -primals} \\
\texttt{idiot60} & \texttt{allInt}, avg $> 0.0086$ &
  \texttt{-idiot 60 -primals} \\
\texttt{dualpesteep} & \texttt{!allInt}, median $\le 0.75$ &
  \texttt{-dualpivot pesteep -psi 1.0 -pertv 52 -duals} \\
\texttt{idiot80} & \texttt{!allInt}, median $> 0.75$ &
  \texttt{-idiot 80 -primals} \\
\bottomrule
\end{tabular}
\end{table}

\subsection{Protocol}
Every instance was solved by \emph{all four} configurations independently (i.e.,
ignoring the actual \texttt{guess()} decision), to enable direct head-to-head
comparison. Per-instance time limits were taken from
\texttt{suggested\_lp\_time\_limits.tsv} (formula:
$\min(7200,\max(2t_{\text{ref}},\, t_{\text{ref}}+1800))$ where $t_{\text{ref}}$
is the \texttt{initialSolve} time on the same instance). Jobs were run in
parallel with 120 cores using GNU parallel. Total jobs: $1317 \times 4 = 5268$.

% ─────────────────────────────────────────────────────────────────────────────
\section{Correctness and Reliability}

\subsection{Memory Safety}
\textbf{Zero ASan errors} were detected across all 5\,268 runs. No assertion
failures were observed either. All four configurations are memory-safe on this
benchmark set.

\subsection{Per-Method Problem Summary}

Table~\ref{tab:problems} summarises correctness and reliability issues per method.
A ``problem'' is any of: killed by the time limit (no solution), objective not
parseable (\texttt{NA}), or objective value differing from the Gurobi reference
by more than a relative tolerance of $10^{-6}$.

\begin{table}[h]
\centering
\caption{Correctness and reliability per configuration (out of 1\,317 instances).}
\label{tab:problems}
\begin{tabular}{lrrrrrr}
\toprule
\textbf{Method} & \textbf{Solved} & \textbf{Killed/NA} & \textbf{False} &
\textbf{Total problems} & \textbf{Problem rate} \\
\midrule
\texttt{idiot30}     & 1276 & 40 & 1 & 41 & 3.1\% \\
\texttt{idiot60}     & 1284 & 32 & 1 & 33 & 2.5\% \\
\texttt{idiot80}     & 1283 & 33 & 1 & 34 & 2.6\% \\
\texttt{dualpesteep} & \best{1285} & \best{30} & \bad{2} & \best{32} & \best{2.4\%} \\
\bottomrule
\end{tabular}
\end{table}

\noindent \textbf{Key observations:}
\begin{itemize}
  \item \texttt{idiot30} has the highest problem rate, driven by the most time-limit
        kills (40). The idiot algorithm with only 30 iterations is the least reliable
        on harder instances.
  \item \texttt{dualpesteep} has the fewest kills (30) but the most false objective
        matches (2), indicating an accuracy issue in specific cases (see
        Section~\ref{sec:false}).
  \item \texttt{idiot60} and \texttt{idiot80} are nearly identical in reliability.
\end{itemize}

\subsection{False Objective Matches}
\label{sec:false}

Five runs produced an objective value that does not match the Gurobi reference.
All involve instances where CBC reports \texttt{Unscaled problem has primal
infeasibilities} after declaring optimality, indicating a post-solve scaling
issue that corrupts the reported objective.

\begin{table}[h]
\centering
\caption{False objective matches (reference value from Gurobi).}
\begin{tabular}{llrrl}
\toprule
\textbf{Instance} & \textbf{Method} & \textbf{CBC obj} & \textbf{Gurobi ref} & \textbf{Log message} \\
\midrule
\texttt{gsvm2rl5}         & \texttt{dualpesteep} & 0.19932 & 0.20000 & primal infeasibilities \\
\texttt{irish-electricity}& \texttt{dualpesteep} & 2\,454\,717 & 2\,546\,255 & primal infeas.\ + barrier\\
\texttt{irish-electricity}& \texttt{idiot30}     & 2\,454\,743 & 2\,546\,255 & primal infeas.\ + barrier\\
\texttt{irish-electricity}& \texttt{idiot60}     & 2\,454\,759 & 2\,546\,255 & primal infeas.\ + barrier\\
\texttt{irish-electricity}& \texttt{idiot80}     & 2\,454\,709 & 2\,546\,255 & primal infeas.\ + barrier\\
\bottomrule
\end{tabular}
\end{table}

\noindent The \texttt{irish-electricity} instance fails on \emph{all four} methods,
suggesting a problem-specific pathology (the instance triggers the barrier solver
internally, which then fails to converge to a primal-feasible point). The
\texttt{gsvm2rl5} failure is specific to \texttt{dualpesteep} only.

\subsection{Instances with No Solution (NA / Killed)}
Of the 58 instances that produced at least one NA or kill across any method:
\begin{itemize}
  \item \textbf{9 instances} timed out on all four methods — these are genuinely
        hard LP relaxations within the per-instance time budget.
  \item \textbf{49 instances} failed on only some methods, indicating
        method-specific difficulty. \texttt{idiot30} is disproportionately
        represented in the method-specific failures.
\end{itemize}

% ─────────────────────────────────────────────────────────────────────────────
\section{Performance Analysis}

\subsection{Summary Statistics}

Table~\ref{tab:timing} reports timing statistics over the 1\,257 instances where
all four methods solved correctly within the time limit.

\begin{table}[h]
\centering
\caption{Timing statistics on 1\,257 commonly solved instances (seconds).}
\label{tab:timing}
\begin{tabular}{lrrrr}
\toprule
\textbf{Method} & \textbf{Median} & \textbf{Mean} & \textbf{Geom.\ mean} & \textbf{Max} \\
\midrule
\texttt{idiot30}     & 6.70 & 219.7 & 7.27 & 8238.6 \\
\texttt{idiot60}     & 7.56 & 190.2 & 7.41 & 8915.9 \\
\texttt{idiot80}     & 8.33 & 177.6 & 7.91 & 7226.4 \\
\texttt{dualpesteep} & \best{2.99} & \bad{287.7} & \best{5.07} & \bad{8747.6} \\
\bottomrule
\end{tabular}
\end{table}

\noindent \texttt{dualpesteep} has the lowest median and geometric mean times,
but the highest arithmetic mean — a consequence of catastrophic slowdowns on a
subset of instances (see Section~\ref{sec:extreme}).

\subsection{Fastest-Method Count and Rank Distribution}

\begin{table}[h]
\centering
\caption{How often each method ranks 1st (fastest) and last (slowest) across
1\,257 commonly solved instances.}
\label{tab:ranks}
\begin{tabular}{lrrrr}
\toprule
\textbf{Method} & \textbf{Rank 1 (fastest)} & \textbf{Rank 2} & \textbf{Rank 3} & \textbf{Rank 4 (slowest)} \\
\midrule
\texttt{idiot30}     & 156 (12.4\%) & 537 & 294 & 263 \\
\texttt{idiot60}     & 176 (14.0\%) & 360 & 490 & 225 \\
\texttt{idiot80}     & 152 (12.1\%) & 249 & 384 & \bad{466 (37.1\%)} \\
\texttt{dualpesteep} & \best{770 (61.3\%)} & 102 & 80  & 300 \\
\bottomrule
\end{tabular}
\end{table}

\texttt{dualpesteep} is the fastest method on 61\% of instances. \texttt{idiot80}
is the slowest most often (37\% of instances), making it the weakest of the four
configurations.

\subsection{Performance Profile}

Table~\ref{tab:profile} shows the fraction of instances each method solves within
a factor $\tau$ of the best (fastest) method on that instance — the standard
performance profile metric~\cite{dolan2002}.

\begin{table}[h]
\centering
\caption{Performance profile: fraction of instances within $\tau \times$ best time.}
\label{tab:profile}
\begin{tabular}{lrrrrrrr}
\toprule
\textbf{Method} & $\tau{=}1.0$ & $\tau{=}1.25$ & $\tau{=}1.5$ & $\tau{=}2$ & $\tau{=}3$ & $\tau{=}5$ & $\tau{=}10$ \\
\midrule
\texttt{idiot30}     & 0.126 & 0.301 & 0.446 & 0.607 & 0.741 & 0.839 & 0.924 \\
\texttt{idiot60}     & 0.140 & 0.350 & 0.474 & 0.611 & 0.738 & 0.847 & 0.918 \\
\texttt{idiot80}     & 0.122 & 0.305 & 0.434 & 0.580 & 0.712 & 0.826 & 0.913 \\
\texttt{dualpesteep} & \best{0.615} & \best{0.713} & \best{0.758} & \best{0.796} & \best{0.850} & \best{0.908} & \best{0.960} \\
\bottomrule
\end{tabular}
\end{table}

\texttt{dualpesteep} dominates the performance profile at every threshold.
At $\tau=1$ it is within the best time on 61.5\% of instances (i.e.\ it \emph{is}
the best method). At $\tau=10$ it is still within $10\times$ best on 96\% of
instances, versus 91--92\% for the idiot variants.

\subsection{Geometric Mean Ratio vs Virtual Best}

A complementary measure is the geometric mean of the ratio $t_{\text{method}} /
t_{\text{best}}$ over all instances:

\begin{center}
\begin{tabular}{lr}
\toprule
\textbf{Method} & \textbf{Geom.\ mean ratio} \\
\midrule
\texttt{idiot30}     & 2.176 \\
\texttt{idiot60}     & 2.142 \\
\texttt{idiot80}     & 2.304 \\
\texttt{dualpesteep} & \best{1.466} \\
\bottomrule
\end{tabular}
\end{center}

\noindent On average, \texttt{dualpesteep} is $1.47\times$ the best method, while
the idiot variants are $2.1$--$2.3\times$ the best.

\subsection{Pairwise Comparison}

Table~\ref{tab:pairwise} reports pairwise median runtime ratios. A ratio $< 1$
means the row method is faster than the column method on the median instance.

\begin{table}[h]
\centering
\caption{Pairwise median runtime ratio (row / column). Values $<1$ indicate
the row method is faster.}
\label{tab:pairwise}
\begin{tabular}{lcccc}
\toprule
 & \texttt{idiot30} & \texttt{idiot60} & \texttt{idiot80} & \texttt{dualpesteep} \\
\midrule
\texttt{idiot30}     & --- & 0.944 & 0.898 & 2.244 \\
\texttt{idiot60}     & 1.059 & --- & 0.929 & 2.393 \\
\texttt{idiot80}     & 1.114 & 1.076 & --- & 2.627 \\
\texttt{dualpesteep} & \best{0.446} & \best{0.418} & \best{0.381} & --- \\
\bottomrule
\end{tabular}
\end{table}

\noindent Among the idiot methods, \texttt{idiot30} is fastest overall (median
ratio $< 1$ vs both idiot60 and idiot80). The differences between the three idiot
variants are modest (5--11\% on the median), while \texttt{dualpesteep} is
roughly $2.2$--$2.6\times$ faster than any idiot variant on the median instance.

\subsection{Significant Wins (Method-Specific Advantage)}

A method has a ``significant win'' on an instance if it is the fastest and more
than $2\times$ faster than the next best method. This identifies instances where a
particular method provides a qualitative advantage.

\begin{center}
\begin{tabular}{lr}
\toprule
\textbf{Method} & \textbf{Significant wins ($>2\times$ faster)} \\
\midrule
\texttt{idiot30}     & 10 \\
\texttt{idiot60}     & 0 \\
\texttt{idiot80}     & 0 \\
\texttt{dualpesteep} & \best{365} \\
\bottomrule
\end{tabular}
\end{center}

\begin{itemize}
  \item \texttt{dualpesteep} dominates: 365 instances with $>2\times$ advantage,
        including extreme cases up to $54\times$ faster (Table~\ref{tab:extreme}).
  \item \texttt{idiot30} has 10 specific wins, concentrated on bin-packing and
        set-covering type instances (e.g.\ \texttt{bppc6-02}, \texttt{bppc6-06},
        \texttt{cl\_06\_020\_*}).
  \item \textbf{\texttt{idiot60} and \texttt{idiot80} have zero significant wins.}
        Neither method is ever more than $2\times$ faster than all alternatives on
        any instance. This strongly suggests that at least one (and possibly both)
        of these configurations could be removed from the \texttt{guess()} selection
        pool without measurable loss.
\end{itemize}

\subsection{Extreme Cases}
\label{sec:extreme}

Table~\ref{tab:extreme} shows the five instances where \texttt{dualpesteep} is
most advantageous and most disadvantageous relative to the best idiot method.

\begin{table}[h]
\centering
\caption{Extreme performance cases for \texttt{dualpesteep}.}
\label{tab:extreme}
\begin{tabular}{llrrl}
\toprule
\textbf{Instance} & \textbf{Type} & \textbf{dualpesteep} & \textbf{Best idiot} & \textbf{Ratio} \\
\midrule
\texttt{neos-5129192-manaia} & win & 17.2s & 920.8s & $53.6\times$ faster \\
\texttt{neos-3755335-nizao}  & win &  4.0s & 134.4s & $33.5\times$ faster \\
\texttt{pb-fit2d}            & win & 28.6s & 921.6s & $32.2\times$ faster \\
\texttt{A-2}                 & win &  0.7s &  21.9s & $30.0\times$ faster \\
\texttt{neos-957270}         & win &  0.8s &  23.7s & $28.6\times$ faster \\
\midrule
\texttt{j12016\_10}          & loss & 6196.4s &  79.4s & $78.0\times$ slower \\
\texttt{supportcase42}       & loss & 7270.5s &  93.8s & $77.5\times$ slower \\
\texttt{scpm1}               & loss & 6791.4s & 191.2s & $35.5\times$ slower \\
\texttt{neos-933966}         & loss &  326.8s &   9.7s & $33.7\times$ slower \\
\texttt{neos-935769}         & loss &  173.0s &   5.5s & $31.6\times$ slower \\
\bottomrule
\end{tabular}
\end{table}

The catastrophic losses of \texttt{dualpesteep} suggest that the \texttt{guess()}
branching condition (based on \texttt{allInt} and coefficient statistics) sometimes
selects this configuration for instances where it performs extremely poorly.

% ─────────────────────────────────────────────────────────────────────────────
\section{Conclusions and Recommendations}

\subsection{Memory Safety}
All four configurations are memory-safe: \textbf{0 ASan errors, 0 assertion
failures} across 5\,268 runs.

\subsection{Correctness}
\begin{itemize}
  \item \texttt{irish-electricity} fails on all four methods due to an internal
        barrier solver convergence issue — this is an instance-level pathology,
        not configuration-specific.
  \item \texttt{gsvm2rl5} fails on \texttt{dualpesteep} only (\texttt{Unscaled
        problem has primal infeasibilities}). The other three methods solve it
        correctly. This is a mild correctness concern for the dual pesteep path.
\end{itemize}

\subsection{Reliability}
\texttt{dualpesteep} has the fewest kills (30 vs.\ 32--40 for idiot methods) and
the lowest overall problem rate (2.4\%). \texttt{idiot30} is the least reliable.

\subsection{Performance}
\begin{itemize}
  \item \texttt{dualpesteep} is by far the strongest method: fastest on 61\% of
        instances, geometric mean ratio $1.47\times$ best vs.\ $2.1$--$2.3\times$
        for idiot methods, and 365 significant wins ($>2\times$ advantage).
  \item Among idiot methods, \texttt{idiot30} is the best overall (lowest
        geometric mean, highest rank-1 count among the three), with 10 specific
        wins on set-covering/bin-packing instances.
  \item \textbf{\texttt{idiot60} and \texttt{idiot80} have no significant wins
        and are nearly indistinguishable from each other.} Removing one or both
        from the \texttt{guess()} portfolio would not sacrifice measurable
        performance on any instance class tested.
\end{itemize}

\subsection{Recommended Action}
\begin{enumerate}
  \item Investigate the \texttt{gsvm2rl5} / \texttt{dualpesteep} infeasibility
        issue — it may be a scaling or perturbation artefact specific to the
        pesteep pivot rule.
  \item Consider replacing \texttt{idiot60}/\texttt{idiot80} in \texttt{guess()}
        with \texttt{idiot30}, retaining only two configurations
        (\texttt{idiot30} + \texttt{dualpesteep}) differentiated by the
        \texttt{allInt} flag.
  \item The 9 instances that time out on all four methods may benefit from a
        fallback to the standard dual simplex (\texttt{initialSolve}).
\end{enumerate}

\begin{thebibliography}{9}
\bibitem{dolan2002}
  E.\ D.\ Dolan and J.\ J.\ Mor\'e,
  \textit{Benchmarking optimization software with performance profiles},
  Mathematical Programming, 91(2):201--213, 2002.
\end{thebibliography}

\end{document}
